\chapter{Conservative Holonomic Systems}
\label{chs}
\section [(In)equivalence of the Formulations]{On the (In)equivalence of the Formulations of Classical Mechanics}
\label{chs:nf}

The Newtonian, Lagrangian, Hamiltonian, and Hamilton--Jacobi formulations of
classical mechanics are routinely presented as equivalent packagings of a
single physical theory, related by Hamilton's principle, the Legendre
transform, and Jacobi's theorem. This claim is correct, but only under
conditions that are rarely stated explicitly and that fail even for
holonomic systems considered above.

Here we show a number  of
 distinct ways in which the standard equivalence breaks down: 
\begin{itemize}
   \item the gap
between Hamilton's principle (stationarity of the action) and the popularly
stated ``principle of \emph{least} action'' (minimality), which fails beyond a
conjugate time fixed by the curvature of the potential and which can be made
arbitrarily short;
   \item the consequent failure of the boundary-value problem built
into Hamilton's principle itself to be well-posed---even for $H=p^2/2m+V(q)$, a
sufficiently steep $V$ makes the two-endpoint problem admit no solution, or
infinitely many, in contrast to the always-well-posed Newtonian initial-value
problem;
   \item the existence of singular Lagrangians, for which the Legendre
transform is not invertible and Dirac's theory of constraints is required;
   \item the existence of Newtonian force laws that admit no
Lagrangian description at all, governed by the Helmholtz conditions;
   \item the
failure of the Legendre transform to be a global diffeomorphism even for
pointwise-regular Lagrangians, absent the stronger condition of
hyperregularity;
   \item the mechanistically identical failure of the Hamilton--Jacobi
equation's generating function to stay single-valued beyond a caustic;
\end{itemize}

Let us display shortly the four standard formulations.

\paragraph{Newtonian mechanics}
For a system of $N$ particles with positions $\bm{r}_i \in \mathbb{R}^3$,
Newton's second law reads $m_i \ddot{\bm{r}}_i = \bm{F}_i$, where
$\bm{F}_i$ is the total force on particle $i$. Constraints are handled by
splitting the force into an applied part and a constraint (reaction) part,
$\bm{F}_i = \bm{F}_i^{\mathrm{app}} + \bm{F}_i^{\mathrm{c}}$, and
postulating---this is d'Alembert's principle---that the constraint forces do no
virtual work: $\sum_i \bm{F}_i^{\mathrm{c}}\cdot \delta \bm{r}_i = 0$
for every virtual displacement $\delta\bm{r}_i$ compatible with the
constraints at a frozen instant of time. It is d'Alembert's principle, not
Newton's second law by itself, that supplies the bridge to the variational
formulations below.

\paragraph{Lagrangian mechanics}
Let $Q$ be the configuration manifold and $TQ$ its tangent bundle, with induced
coordinates $(q^i,\dot q^i)$. A Lagrangian is a function $L : TQ \times
\mathbb{R} \to \mathbb{R}$, typically (but not necessarily) of the form $L = T
- V$, kinetic minus potential energy. Hamilton's principle asserts that
physical trajectories $q(t)$ extremize the action
\begin{equation}
S[q(\cdot)] = \int_{t_0}^{t_1} L\big(q(t),\dot q(t), t\big)\, dt
\end{equation}
among curves with fixed endpoints $q(t_0), q(t_1)$. Setting the first variation
to zero yields the Euler--Lagrange equations
\begin{equation}
\label{EL}
\frac{d}{dt}\frac{\partial L}{\partial\dot q^i} - \frac{\partial L}{\partial q^i} = 0 ,
    \qquad i = 1,\dots,n=\dim Q .
\end{equation}
For a system subject to \emph{ideal, holonomic} constraints, this is fully
equivalent to Newton's second law plus d'Alembert's principle, once $Q$ is
taken to be the constraint surface itself and the $q^i$ are independent
generalized coordinates on it.

\paragraph{Hamiltonian mechanics}
Define the momenta $p_i = \partial L/\partial \dot q^i$. This is the
\emph{fiber derivative} 
$\mathbb{F}L: TQ \to T^*Q$, $(q,\dot q)\mapsto (q,p)$.
Where $\mathbb{F}L$ can be inverted for $\dot q$ in terms of $(q,p)$, one
defines the Hamiltonian by the Legendre transform,
\begin{equation}
H(q,p,t) = p_i \dot q^i - L(q,\dot q, t)\Big|_{\dot q = \dot q(q,p,t)} ,
\end{equation}
and Hamilton's canonical equations \footnote{ These equations also can 
be obtained from variational principle discussed later}
\begin{equation}
\label{Ham}
\dot q^i = \frac{\partial H}{\partial p_i}, \qquad 
    \dot p_i = -\frac{\partial H}{\partial q^i}
\end{equation}
reproduce~\eqref{EL} on the nose. The cotangent bundle $T^*Q$ carries a
canonical symplectic form $\omega = dp_i \wedge dq^i$, giving Hamiltonian
mechanics its distinctive geometry: canonical transformations, Poisson
brackets, Liouville's theorem, and so on.

The precise statement licensing the textbook claim is the following.
\begin{proposition}[Regular equivalence]
\label{regular}
Let $$L: TQ\times\mathbb{R}\to\mathbb{R}$$
    be \emph{regular}: the fiber Hessian
    $$W_{ij} = \partial^2 L/\partial \dot q^i \partial \dot q^j$$
    is nondegenerate at every point of $TQ$. Then $\mathbb{F}L$ is a local
    diffeomorphism, $H$ is well defined locally, and solutions of~\eqref{EL}
    correspond, via 
    $$p = \mathbb{F}L(q,\dot q),$$ 
    bijectively and locally to
    solutions of~\eqref{Ham}. If in addition $\mathbb{F}L$ is a \emph{global}
    diffeomorphism $TQ \to T^*Q$---in which case $L$ is called
    \emph{hyperregular}---this correspondence extends to a global bijection
    between the full solution spaces of the two theories.
\end{proposition}

\paragraph{Hamilton--Jacobi theory}
A fourth formulation, historically inseparable from Hamilton's own work,
replaces the $2n$ first-order equations~\eqref{Ham} by a single first-order
partial differential equation for a scalar function. Define \emph{Hamilton's
principal function} $S(q,t)$ as the action evaluated along the true trajectory
of the system that starts at some fixed reference event $(q_0,t_0)$ and ends at
$(q,t)$. Hamilton showed that, viewed as a function of its endpoint, $S$
satisfies the \emph{Hamilton--Jacobi equation}
\begin{equation}
\label{HJ}
\frac{\partial S}{\partial t}+H\Big(q,\frac{\partial S}{\partial q},t\Big)=0.
\end{equation}
Conversely, Jacobi showed how to run this construction backwards and recover
the canonical equations~\eqref{Ham} directly from~\eqref{HJ}, without
integrating an ordinary differential equation, provided one can find a
\emph{complete integral}: an $n$-parameter family of solutions $S(q,t;\alpha)$,
$\alpha=(\alpha_1,\dots,\alpha_n)$, of~\eqref{HJ} whose mixed Hessian
$\partial^2 S/\partial q^i\partial \alpha_j$ is everywhere nondegenerate on the
\emph{domain of interest}.

\begin{proposition}[Jacobi's theorem]
\label{jacobi-thm}
Let $S(q,t;\alpha)$ be a complete integral of~\eqref{HJ} on an open set
    $U\subset Q\times\mathbb{R}$, with $\det(\partial^2 S/\partial
    q^i\partial\alpha_j)\neq 0$ throughout $U$. Then the equations
\begin{equation}
\frac{\partial S}{\partial\alpha_j} = \beta_j \ (\text{constants}), 
    \qquad p_i = \frac{\partial S}{\partial q^i}
\end{equation}
implicitly define, for each choice of constants $(\alpha,\beta)$, a curve
    $(q(t),p(t))$ solving Hamilton's equations~\eqref{Ham} on $U$; every
    solution of~\eqref{Ham} passing through $U$ arises this way for a suitable
    choice of $(\alpha,\beta)$.
\end{proposition}

Proposition~\ref{jacobi-thm} is the formal engine behind solving
Hamiltonian systems by separation of variables, and it is a genuine equivalence
result: wherever a complete integral with nondegenerate mixed Hessian exists,
the Hamilton--Jacobi and Hamiltonian pictures determine exactly the same
trajectories. The qualification \emph{'on the domain of interest'} is doing real
work, and  cannot in general be
dropped, and that its failure is tied to the same mechanism that undermines the
``least action'' reading of Hamilton's principle.

Everything that follows  concerns ways in which the
hypotheses of Proposition~\ref{regular} or
Proposition~\ref{jacobi-thm} can fail to hold, or ways in which their
being satisfied does not settle every reasonable sense in which one might want
two physical theories to be ``the same.'' 

The equivalence between Hamilton's principle and the Euler--Lagrange
equations \eqref{EL} is exact and unconditional: $\delta S=0$ for all
admissible variations if and only if \eqref{EL} holds along the trajectory.
This is a first-order, local statement about the action, and none of
standard formulations depends on anything more. But the popular name for
Hamilton's principle---the ``principle of \emph{least} action''---asserts
something stronger: that the true trajectory does not merely make $S$
stationary but \emph{minimizes} it among nearby competitors with the same
endpoints. This stronger claim concerns the second variation, not the first,
and it is generally false, even for the most elementary holonomic,
unconstrained, regular Lagrangian system one could write down.

For a one-dimensional natural Lagrangian 
$$L(x,\dot x) = \tfrac12 m\dot x^2 - V(x),$$
consider a perturbation $x(t)+\varepsilon\eta(t)$ of a solution $x(t)$
of \eqref{EL}, with $\eta(t_0)=\eta(t_1)=0$. The second variation of the action
is
\begin{equation}
\delta^2 S[\eta] = \int_{t_0}^{t_1} \Big[ m\dot\eta(t)^2 - V''(x(t))\,\eta(t)^2 \Big]\, dt ,
\end{equation}
and the trajectory is a genuine local minimum of $S$ if and only if
$\delta^2S[\eta]>0$ for every admissible $\eta\not\equiv0$. Whether this holds
is governed by the \emph{Jacobi accessory equation}, obtained by
linearizing~\eqref{EL} about $x(t)$,
\begin{equation}
\label{jacobi-eq}
m\ddot\eta + V''(x(t))\,\eta = 0 ,
\end{equation}
together with $\eta(t_0)=0$. A time $t_c>t_0$ is a \emph{conjugate point} (or
\emph{kinetic focus}) to $t_0$ along the trajectory if the nontrivial solution
of~\eqref{jacobi-eq} with $\eta(t_0)=0$ vanishes again at $t_c$.

\begin{proposition}[Jacobi's necessary condition]
\label{jacobi-necessary}
The trajectory $x(t)$ furnishes a local minimum of the action on $[t_0,t_1]$
    only if $t_1$ occurs before the first conjugate point to $t_0$. If $t_1$
    lies beyond the first conjugate point, $S$ is a saddle point at $x(t)$:
    nearby trial trajectories exist with strictly smaller action, even though
    $x(t)$ exactly satisfies~\eqref{EL}.
\end{proposition}

\begin{example}[Harmonic oscillator]
\label{sho}
Take $V(x)=\tfrac12 kx^2$, so $V''\equiv k$ and \eqref{jacobi-eq} becomes
    $\ddot\eta+\omega^2\eta=0$ with $\omega=\sqrt{k/m}$. The solution with
    $\eta(t_0)=0$ is $\eta(t)=\sin\!\big(\omega(t-t_0)\big)$, which next
    vanishes at $t_c=t_0+\pi/\omega$---exactly half a period. By
    Proposition~\ref{jacobi-necessary}, a segment of oscillator trajectory
    genuinely minimizes the action if its duration is less than $\pi/\omega$,
    and is merely a stationary saddle point if its duration exceeds
    $\pi/\omega$.
\end{example}

Example~\ref{sho} answers, quantitatively, the question of how short is ``short
enough.'' The conjugate-point time $\pi/\omega=\pi\sqrt{m/k}$ is set entirely
by the curvature $V''=k$ of the potential and shrinks without bound as $k$
grows: no fixed time interval is safe for every system, because a sufficiently
steeply curved potential will already have produced a conjugate point within
it. More generally, conjugate points arise only where $V''>0$ along the
trajectory; for potentials with $V''(x)\leq 0$ for all $x$ (a free particle,
uniform gravity, an inverted quadratic barrier) no conjugate point ever forms,
and the action is minimized for arbitrarily long times.
Whether ``the action is minimized'' is therefore not a property of Hamilton's
principle in general---it is a property of the specific potential, and
specifically of how sharply it curves.

It is worth being precise about what this does and does not show. Hamilton's
principle, in its correct form ($\delta S=0$), remains exactly equivalent to
Newton's second law throughout; nothing here contradicts
Proposition~\ref{regular}. What fails is the stronger, popularly stated
identification of Hamilton's principle with a principle of \emph{least}
action---and it fails already in the cleanest setting available anywhere in
this paper: a single unconstrained particle, regular and hyperregular
Lagrangian, smooth potential. That is precisely why this failure mode belongs
at the head of this subsection, ahead of the more exotic constructions that
follow.

The question is asked whether a trajectory already known to
satisfy~\eqref{EL} minimizes the action. A logically prior question is whether
such a trajectory exists and is unique at all, given the data Hamilton's
principle is actually stated in terms of: not an initial position and velocity,
but two fixed \emph{endpoints}, $q(t_0)=q_A$ and $q(t_1)=q_B$---exactly the
boundary conditions built into the action functional. 
This is a fundamentally different mathematical
problem from the Newtonian initial-value problem. Given $(q_0,\dot q_0)$ at a
single instant, standard ODE theory (Picard--Lindelöf) guarantees a unique
solution of~\eqref{EL}, at least locally in time, for any reasonably smooth
force law: the ``most boring'' Hamiltonian, $H=p^2/2m+V(q)$, is always
well-posed as an initial-value problem. Fixing $q(t_0)$ and $q(t_1)$ instead
offers no such guarantee. The boundary-value problem can have no solution, a
unique solution, or infinitely many, and which of these holds is governed by
exactly the conjugate points.

Fix $t_0$ and $q_A=q(t_0)$, and let $x(t,v_0)$ denote the unique initial-value
solution with $\dot x(t_0)=v_0$. The \emph{shooting map} $v_0\mapsto
x(t_1,v_0)$ determines which endpoints are reachable in time $T=t_1-t_0$ and
how many initial velocities reach a given one. Its derivative $\partial
x(t,v_0)/\partial v_0$ satisfies~\eqref{EL} linearized about the
trajectory---for $H=p^2/2m+V(q)$ this is exactly the Jacobi accessory
equation~\eqref{jacobi-eq}, with initial data $\eta(t_0)=0$, $\dot\eta(t_0)=1$.
Consequently
\begin{equation}
\frac{\partial x(t_1,v_0)}{\partial v_0} = 0
    \quad\Longleftrightarrow\quad t_1 \text{ is a conjugate point to } t_0 ,
\end{equation}
so the shooting map is a local diffeomorphism precisely up to the first
conjugate time, and can fold---losing injectivity, surjectivity, or both---at
and beyond it.

\begin{example}[Harmonic oscillator boundary-value problem]
\label{sho-bvp}
Continuing Example~\ref{sho}, set $t_0=0$, $q_A=0$. The initial-value solution
    is $x(t,v_0)=(v_0/\omega)\sin(\omega t)$, so the shooting map to time $T$
    is $v_0\mapsto (v_0/\omega)\sin(\omega T)$, and
\begin{equation}
v_0 = \frac{\omega\, q_B}{\sin(\omega T)}
\end{equation}
solves $x(0)=0$, $x(T)=q_B$ whenever $\sin(\omega T)\neq0$: a unique solution,
    as expected for $T$ short of the first conjugate time $\pi/\omega$. But at
    $T=n\pi/\omega$ for a positive integer $n$---an integer multiple of the
    conjugate time---$\sin(\omega T)=0$ identically, and the shooting map
    collapses to $v_0\mapsto 0$ regardless of $v_0$: for $q_B\neq0$ \emph{no}
    solution exists, while for $q_B=0$ \emph{every} $v_0\in\mathbb{R}$ solves
    the boundary-value problem, a one-parameter family of distinct true
    trajectories connecting the same two events.
\end{example}

Example~\ref{sho-bvp} shows non-existence and non-uniqueness sitting on
opposite sides of the same coincidence, both triggered by the same
conjugate-time condition, and both occurring for the most elementary confining
potential there is. The critical travel times $n\pi/\omega$ are evenly spaced
only because the harmonic oscillator's own equation of motion is already
linear---so it coincides with its own Jacobi equation, and conjugate points
recur exactly periodically. For a general potential, the accessory
equation~\eqref{jacobi-eq} has a time-varying coefficient $V''(x(t))/m$ and the
conjugate times need not be evenly spaced, but the qualitative picture
survives, and can be more dramatic still: so quartic
oscillator, $V(x)=\tfrac14 Cx^4$, in which a whole family of distinct true
trajectories leaving a common starting event recross one another at a common
later event, producing multiple genuinely different---not merely differently
scaled---solutions of the same boundary-value problem.

 Since $\omega=\sqrt{k/m}$ for the harmonic case,
the critical travel times $n\pi/\omega = n\pi\sqrt{m/k}$ shrink as the
curvature $k$ grows, so for \emph{any} fixed travel time $T$, a steep enough
potential well places $T$ at or past a degenerate configuration. ``The most
boring system there is,'' $H=p^2/2m+V(q)$, therefore fails to have a well-posed
boundary-value problem in exactly the regime---strongly confining, rapidly
growing $V$---that most textbook potentials are drawn from.

This multiplicity is not a curiosity confined to classical mechanics: it is
exactly what the Feynman path-integral formulation of quantum mechanics must
sum over. The semiclassical propagator between $q_A$ and $q_B$ is, in general,
a sum of contributions from \emph{every} classical trajectory solving the
boundary-value problem, each carrying a phase determined by its action together
with a further phase set by the \emph{Maslov index}---the number of conjugate
points the trajectory has crossed. The Maslov index is, in
this sense, a bookkeeping device  it is needed because the classical
boundary-value problem it
quantizes is not, in general, uniquely solvable.

A Lagrangian is \emph{singular} (or degenerate) when its fiber Hessian $W_{ij}
= \partial^2 L/\partial \dot q^i \partial \dot q^j$ is degenerate somewhere on
$TQ$. In that case $\mathbb{F}L$ is not a local diffeomorphism there: the
momenta $p_i = \partial L/\partial \dot q^i$ are not independent functions of
the velocities, and one cannot solve for all of $\dot q$ in terms of $(q,p)$.
Concretely, degeneracy of $W_{ij}$ means the map $\mathbb{F}L$ satisfies one or
more relations
\begin{equation}
\phi_a(q,p) = 0, \qquad a = 1,\dots,k,
\end{equation}
purely among the momenta and coordinates---\emph{primary constraints}, in
Dirac's terminology---that hold identically on the image of $\mathbb{F}L$,
before any equations of motion are imposed.

The cleanest illustration is to note that Hamiltonian mechanics is itself,
formally, a maximally singular special case of Lagrangian mechanics. Treat
$(q,p)$ as $2n$ independent coordinates on an auxiliary configuration space and
consider the first-order ``Hamiltonian-form'' Lagrangian
\begin{equation}
L_1(q,p,\dot q,\dot p) = p_i \dot q^i - H(q,p).
\end{equation}
Its Euler--Lagrange equations reproduce~\eqref{Ham} exactly, so in one sense
Hamiltonian mechanics \emph{is} a special case of Lagrangian mechanics on a
bigger space. But the fiber Hessian of $L_1$ with respect to $(\dot q,\dot p)$
vanishes identically---$L_1$ is linear, not quadratic, in the velocities---so
$L_1$ is about as singular as a Lagrangian can be, and the naive
Legendre-transform machinery simply does not
apply to it. Recovering Hamilton's equations from $L_1$ requires running the
full constrained-Hamiltonian algorithm below, not the regular Legendre
transform.

The general treatment of singular Lagrangians is due to Dirac and Bergmann.
 One passes to the \emph{primary constraint surface}
$\Sigma_1 \subset T^*Q$ cut out by the $\phi_a$, and checks whether the primary
constraints are preserved under time evolution generated by a (now only
partially determined) Hamiltonian; if not, this consistency requirement
generates \emph{secondary constraints}, and the process is iterated until a
stable constraint surface is reached. Constraints are further classified as
\emph{first class} (those with vanishing Poisson brackets with all other
constraints, which typically signal gauge freedom---directions in phase space
that are not physically distinguishable) or \emph{second class} (which can be
eliminated using the Dirac bracket, an amended Poisson bracket adapted to the
constraint surface). The upshot is that for a
singular $L$, the physical phase space is not $T^*Q$ but a further-reduced
submanifold, of dimension strictly smaller than $2n$ in the presence of
first-class constraints, and the map from Lagrangian solutions to Hamiltonian
solutions is not the clean bijection of Proposition~\ref{regular} but a
considerably more elaborate correspondence, often carrying residual gauge
freedom on the Lagrangian side that has no counterpart in a naively
written-down Hamiltonian system.

Singular Lagrangians are not a pathology confined to toy examples: they are the
generic situation for gauge theories. Electromagnetism, Yang--Mills theory, and
general relativity all have singular Lagrangians, precisely because the local
gauge symmetry (arising from Noether's second theorem) forces degeneracy of the
fiber Hessian. For essentially all of fundamental field theory, then, the
``equivalence'' between the Lagrangian and Hamiltonian pictures is not the
elementary Legendre-transform story at all, but
the considerably more delicate Dirac--Bergmann correspondence.

\paragraph{The inverse problem: not every Newtonian system is Lagrangian}

Even in the unconstrained case, the passage from Newtonian to Lagrangian
mechanics is not guaranteed to exist at all. Given a system of second-order
ODEs $\ddot q^i = f^i(q,\dot q,t)$---an arbitrary Newtonian force law, divided
by mass---one can ask whether it arises as the Euler--Lagrange equations of
\emph{some} Lagrangian. This is the \emph{inverse problem of the calculus of
variations}. The answer is governed by the \emph{Helmholtz conditions}: the
linearized (Fr\'echet) differential operator associated with the equations of
motion must be self-adjoint. Douglas (1941) gave a complete classification
of when a Lagrangian exists in the two-degree-of-freedom case, building on
earlier work of Hirsch and Mayer; the general $n$-dimensional theory remains an
active area of research in the geometric calculus of variations.

The familiar case that fails the Helmholtz conditions is linear
velocity-dependent damping. The damped harmonic oscillator,
\begin{equation}
m\ddot x + \gamma \dot x + k x = 0 ,
\end{equation}
is not the Euler--Lagrange equation of any Lagrangian of the ordinary
autonomous, single-particle form $L(x,\dot x)$: the friction term breaks the
self-adjointness that the Helmholtz conditions require. This is not merely a
failure of imagination; it reflects the fact that a Hamilton's-principle
Lagrangian for a single degree of freedom automatically conserves a symplectic
volume (Liouville's theorem), while a damped oscillator's phase-space volume
visibly contracts as the system loses energy to friction. One can restore a
variational description, but only at a structural cost: by doubling the degrees
of freedom (the Bateman construction, pairing the damped oscillator with a
time-reversed, amplified partner), by using an explicitly time-dependent
Lagrangian such as the Caldirola--Kanai form $L = e^{\gamma t/m}\big(\tfrac12
m\dot x^2 - \tfrac12 kx^2\big)$, or by admitting non-standard
(non-quadratic-in-velocity) Lagrangians. Each of these routes reproduces the
correct equation of motion, but none of them is the ``obvious,''
structure-preserving Lagrangian one would have guessed, and none is unique.

The general moral is that Newtonian mechanics, understood as ``any second-order
ODE system one likes,'' is strictly more permissive than Lagrangian mechanics:
the class of Lagrangian-representable force laws is a proper,
Helmholtz-constrained subset of the class of Newtonian force laws, unless one
is willing to pay for a Lagrangian description with auxiliary variables or
explicit time dependence.

\paragraph{Global obstructions: regularity is not hyperregularity}
Proposition~\ref{regular} distinguishes \emph{regularity} (the fiber
Hessian is nondegenerate at each point) from \emph{hyperregularity} (the fiber
derivative $\mathbb{F}L$ is a diffeomorphism of the whole of $TQ$ onto the
whole of $T^*Q$). The distinction matters because regularity is a purely local,
pointwise condition, while the equivalence of the full \emph{solution spaces}
of the Lagrangian and Hamiltonian theories is a global statement about
$\mathbb{F}L$.

A Lagrangian can be regular everywhere and still fail to be hyperregular, in
either of two ways. First, $\mathbb{F}L$ can fail to be \emph{surjective}: if
$L$ grows sub-quadratically in $\dot q$ at large velocity (rather than the
quadratic-plus-lower-order growth of the usual kinetic-minus-potential form),
the momenta $p=\partial L/\partial\dot q$ can be bounded even as $\dot
q\to\infty$, so that some momentum values in $T^*Q$ are simply never reached,
and the corresponding Hamiltonian trajectories have no Lagrangian counterpart.
Second, $\mathbb{F}L$ can fail to be \emph{injective} globally even while being
a local diffeomorphism everywhere---the fiber derivative can ``fold back'' on
itself far from any particular point, so that two distinct velocities at the
same base point map to the same momentum. In either case, one is forced either
to restrict attention to a proper open subset of phase space on which
$\mathbb{F}L$ does behave well, or to accept that the Lagrangian and
Hamiltonian pictures describe genuinely different (if overlapping) collections
of physically possible histories. This is a comparatively mild failure mode,
 but it shows that even the
textbook regularity condition, usually checked (if at all) by verifying $\det
W_{ij}\neq 0$ at a single representative point, is not by itself sufficient for
the global equivalence the folklore claims.

\paragraph{The Hamilton--Jacobi correspondence: caustics}
Jacobi's theorem (Proposition~\ref{jacobi-thm}) is, like
Proposition~\ref{regular}, a genuine equivalence result---but, like
Proposition~\ref{regular} without hyperregularity, it is only a
\emph{local} one, and the nondegeneracy hypothesis on the mixed Hessian
$\partial^2S/\partial q\partial\alpha$ is not guaranteed to persist away from
the point where it is first checked.

Construct Hamilton's principal function concretely,
fix an initial event $(q_0,t_0)$, and for each
nearby $(q,t)$ let $S(q,t)$ be the action along the trajectory that solves
Hamilton's equations with $q(t_0)=q_0$ and $q(t)=q$. For $S$ to be a
well-defined, single-valued, smooth function of $(q,t)$, this connecting
trajectory must exist and be \emph{unique}. Existence is guaranteed for short
enough time by ordinary differential equation theory. Uniqueness is exactly
where the conjugate points re-enter: consider the
family of trajectories leaving $q_0$ at $t_0$ with varying initial velocity,
and ask where two members of this family, with infinitesimally different
initial velocity, next intersect. This is precisely the kinetic-focus
construction, and it recurs at the same time $t_c$
computed there. Beyond $t_c$, distinct nearby trajectories from $(q_0,t_0)$ can
reconverge near the same later point, so more than one candidate value competes
to be $S(q,t)$: the function becomes multivalued, its would-be smooth branches
meeting along an envelope---the \emph{caustic}---where the Jacobian of the map
from initial velocity to endpoint, and with it the mixed Hessian of
Proposition~\ref{jacobi-thm}, degenerates.

\begin{proposition}
\label{caustic}
Let $t_c$ be the first conjugate time to $t_0$ along a trajectory.
    Hamilton's principal function $S(q,t)$, built
    from trajectories issuing from $(q_0,t_0)$, is a smooth, single-valued
    solution of the Hamilton--Jacobi equation~\eqref{HJ} for $t<t_c$, and
    generically fails to remain single-valued for $t>t_c$, where the family of
    trajectories develops a caustic.
\end{proposition}

Nothing dynamically singular happens to the particle at a caustic: the actual
trajectories---the characteristics of~\eqref{HJ}---pass through it perfectly
smoothly, exactly as they pass through the conjugate point.
What breaks down is the specific representational
device of Hamilton--Jacobi theory: a single, smooth, globally valid generating
function from which the entire flow can be read off via $p=\partial S/\partial
q$. This is the same phenomenon, and the same underlying mathematics, that
produces caustics in geometric optics---where Hamilton first developed the
parallel theory---and it is the classical root of the phase corrections (the
Maslov index) that enter semiclassical, WKB-type quantization once a trajectory
has passed through one or more foci.

``The Hamiltonian and
Hamilton--Jacobi formulations are equivalent'' is true, without qualification,
only up to the first conjugate point; beyond it, recovering the dynamics from
Hamilton--Jacobi theory requires patching together multiple local branches of
$S$, or abandoning the single smooth generating function in favor of directly
integrating~\eqref{Ham}. This has nothing to do
with singular Lagrangians, nonholonomic constraints, or a failure of
hyperregularity: the same one-dimensional particle in a smooth potential well
illustrates this failure too, because
a conjugate point in the variational sense and a point on a caustic in the
Hamilton--Jacobi sense are the same event, described in two different
formulations of classical mechanics.

\paragraph{Non-uniqueness: the correspondence is not one-to-one in the other direction}
Even restricting entirely to the well-behaved case---holonomic constraints, $L$
regular and hyperregular---the map from ``the physics'' to ``a Lagrangian (or
Hamiltonian) representing it'' is many-to-one rather than one-to-one, which
complicates any claim that a given Lagrangian and a given Hamiltonian
formulation are simply \emph{the} representations of one theory. Adding a total
time derivative, $L \mapsto L + \frac{d}{dt}F(q,t)$ for arbitrary $F$, leaves
the Euler--Lagrange equations, and hence all physical predictions, unchanged;
more drastically, so-called $s$-equivalent Lagrangians can differ by more than
a total derivative and still generate the same force law, again governed by the
multiplier version of the Helmholtz analysis. On
the Hamiltonian side, canonical transformations relate infinitely many
distinct-looking triples $(q,p,H)$ that all describe the same underlying
dynamics. None of this contradicts Proposition~\ref{regular}---the
\emph{solution sets} still correspond bijectively---but it means the intuitive
picture of ``one Newtonian system, one Lagrangian, one Hamiltonian''
understates how much freedom is built into the correspondence in both
directions.

\section[Geometry vs Coordinates]{Geometric Elegance versus Coordinate Practice }
\label{chs:geom}

There are two ways to write down classical mechanics, and every
physicist learns them in the wrong order. First comes the practical
apparatus: pick generalized coordinates $q^i(t)$, write a Lagrangian
$L(q,\dot q,t)$ or a Hamiltonian $H(q,p,t)$, grind through the
Euler--Lagrange or Hamilton equations, solve. This is what one
actually computes with, and it is indispensable --- but it is also,
in a precise sense, an arbitrary choice of language: nothing about
the physics singles out these particular coordinates $q^i$ over some
other set. Only later, if at all, does a student meet the second
formulation: configuration space as an abstract smooth manifold,
velocities and momenta as elements of its tangent and cotangent
bundles, the dynamics generated by a connection or a symplectic
form defined without reference to any coordinate system whatsoever.

The two formulations describe exactly the same physical content; the
tension between them is not a tension between two theories but
between two languages for one theory. And as with any well-chosen
change of language, some things become obvious that were obscure,
and some things become laborious that were once immediate. This
chapter sets the two side by side, with a single worked example run
in both, and asks what each buys --- and costs.

\subsection{The Coordinate Form, and What It Hides}

The familiar machinery needs only a brief reminder. Generalized
coordinates $q^1,\dots,q^n$ parametrize the configurations of a
system; the Lagrangian $L(q,\dot q,t) = T - V$ is a function of
coordinates, velocities, and time; the equations of motion are
\begin{equation}
  \frac{d}{dt}\frac{\partial L}{\partial \dot q^i}
  - \frac{\partial L}{\partial q^i} \;=\; 0,
  \qquad i=1,\dots,n.
  \label{el-coord}
\end{equation}
This is the form in which mechanics is actually solved: one grinds
out~\eqref{el-coord} in whatever coordinates the problem hands
you, and out comes a trajectory.

Consider the simplest possible case: a free particle of mass $m$
moving in a plane, described in plane polar coordinates $(r,\phi)$
rather than Cartesian ones. The kinetic energy is $T =
\tfrac12 m(\dot r^2 + r^2\dot\phi^2)$, and there is no potential, so
$L=T$. Equation~\eqref{el-coord} gives
\begin{equation}
  m\ddot r - m r\dot\phi^2 \;=\; 0,
  \qquad
  \frac{d}{dt}\bigl(m r^2\dot\phi\bigr) \;=\; 0.
  \label{polar-free}
\end{equation}
The second equation is unremarkable --- it just says angular
momentum about the origin is conserved, as it must be for a free
particle. The first is the interesting one: it contains a term
$-mr\dot\phi^2$ that looks, for all the world, like an outward
centrifugal force. There is, of course, no force of any kind acting
on this particle; it travels in a perfectly straight line. The term
is not physics. It is the price of insisting on writing Newton's
second law in a coordinate system whose coordinate lines are
curved --- exactly analogous to the centrifugal and Coriolis terms
that appear when the same free-particle problem is worked in a
rotating frame, except that here the "frame" is not rotating at all;
only the \emph{coordinate grid} is curved.

This is not a quirk of polar coordinates specifically. In any
curvilinear coordinate system, extra terms of exactly this kind are
forced to appear, because the bare second derivative $\ddot q^i$ is
not, by itself, a coordinate-independent (tensorial) object; only the
combination
\begin{equation}
  \ddot q^i + \Gamma^i_{jk}(q)\,\dot q^j\dot q^k
  \label{covariant-accel}
\end{equation}
transforms correctly under a change of coordinates, where $\Gamma^i_{jk}$ are
the Christoffel symbols of the coordinate system (built from the kinetic-energy
metric $ds^2 = \sum_i m_i\,dq^i dq^i$). The Euler--Lagrange machinery
automatically assembles the correctly covariant
combination~\eqref{covariant-accel}, but it does so by generating, coordinate
by coordinate, extra terms that read exactly like additional applied forces ---
centrifugal, Coriolis, Euler --- even when, as here, the true physics is
completely force-free.

None of this is a defect of the coordinate approach; it is the
coordinate approach doing exactly what it is for. This is how
problems actually get integrated, in closed form or numerically, and
it is the only form in which mechanics ever makes contact with a
laboratory: no experiment outputs "a point of an abstract manifold,"
it outputs a number, which is to say a coordinate, read off a ruler,
a protractor, or a clock. A great deal of the technical skill of
mechanics --- finding action-angle variables for an integrable
system, Delaunay or Poincar\'e elements in celestial mechanics,
normal coordinates near an equilibrium --- consists precisely in
finding the coordinates in which a given problem's equations of
motion become as simple as possible. Coordinate choice is not
bookkeeping; it is frequently the whole content of solving a hard
problem.

But this leaves an uncomfortable question hanging. The same physical
system, worked in different coordinates, can produce equations of
motion that look wildly different in complexity, that display or
conceal conserved quantities, that contain "force" terms present in
one description and absent in another. Precisely because the
coordinate form does not, by itself, distinguish physics from
artifact, it cannot by itself answer a question that matters a great
deal: how much of what is written on the page in
equation~\eqref{polar-free} is actually about the world, and how
much is only about the choice of $(r,\phi)$?

\subsection{The Geometric Reformulation}

\paragraph{Configuration space as a manifold}
The geometric answer to that question begins by refusing to
privilege any coordinate system in the first place. Configuration
space $Q$ is taken to be an abstract smooth manifold; a specific
choice of generalized coordinates $q^i$ is only ever a local chart on
$Q$, one among infinitely many equally valid alternatives, and
physical content is, by definition, whatever does not depend on
which chart one happens to be using.

Velocities of the system live in the tangent bundle $TQ$; the Lagrangian is
intrinsically a function $L : TQ\times\mathbb R \to \mathbb R$, stated without
reference to any coordinates at all. The Euler--Lagrange
equations~\eqref{el-coord} themselves have a coordinate-free meaning: they are
the flow of a specific vector field on $TQ$, picked out by the
Poincar\'e--Cartan structure that $L$ induces on $TQ$, satisfying the
second-order condition that this flow really does generate velocities
consistent with positions. What one writes down in coordinates
as~\eqref{el-coord} is the local expression of this single intrinsic object,
valid in every chart, looking different in each.

\paragraph{Phase space and the symplectic form}
Passing (via the Legendre transform, where it is well defined) to
the momenta $p_i = \partial L/\partial \dot q^i$ carries the
description to the cotangent bundle $T^*Q$, the phase space. This
space carries a canonical structure that exists with reference to no
coordinates whatsoever: the \emph{tautological} (or Liouville) 1-form
$\theta$, defined at a point $(q,p)\in T^*Q$ by $\theta(v) =
p\bigl(\pi_*v\bigr)$ for any tangent vector $v$ there, where $\pi:
T^*Q\to Q$ is the projection. Its exterior derivative $\omega =
d\theta$ is the canonical symplectic form on $T^*Q$: closed,
nondegenerate, and built entirely out of the bundle structure of
$T^*Q$ over $Q$, with no reference to a coordinate system anywhere in
its definition.

Hamilton's equations are then the local expression of a single
intrinsic equation,
\begin{equation}
  \iota_{X_H}\,\omega \;=\; dH,
  \label{hamiltonian-vf}
\end{equation}
which defines the Hamiltonian vector field $X_H$ for any smooth
function $H$ on phase space. In local coordinates $(q^i,p_i)$ in
which $\theta = p_i\,dq^i$ and $\omega$ takes its standard constant
form, \eqref{hamiltonian-vf} reproduces exactly the familiar
\begin{equation}
  \dot q^i \;=\; \frac{\partial H}{\partial p_i},
  \qquad
  \dot p_i \;=\; -\,\frac{\partial H}{\partial q^i}.
  \label{hamilton-coord}
\end{equation}
Coordinates in which $\omega$ takes this simple constant form are
called \emph{canonical}, or \emph{Darboux}, coordinates, after Gaston
Darboux, who proved in 1882 that such coordinates always exist, at
least locally, around every point of every symplectic
manifold. This is a striking fact with no analogue
in Riemannian geometry, where curvature is a genuine local invariant
obstructing the metric from being brought to a standard flat form
even at a single point's neighborhood. Symplectic manifolds have no
such local invariant: every symplectic manifold of a given dimension
looks, locally, exactly like every other one. This is precisely why
Hamilton's equations can always be brought to the same simple
form~\eqref{hamilton-coord}, no matter how complicated the
physical system --- the complexity of a hard mechanics problem lives
entirely in the choice of \emph{which} canonical coordinates to use,
never in whether they exist.

\paragraph{Canonical transformations reinterpreted}
In the coordinate treatment, canonical transformations tend to arrive
as a slightly mysterious piece of technology: a table of generating
functions $F_1,F_2,F_3,F_4$, and a set of conditions to check that a
proposed change of variables $(q,p)\mapsto(Q,P)$ preserves the form
of Hamilton's equations. From the geometric point of view, a
canonical transformation is simply a diffeomorphism of phase space
that preserves $\omega$ --- a \emph{symplectomorphism}. The
generating-function apparatus becomes a practical tool for
constructing particular symplectomorphisms, rather than the
definition of the concept. Darboux's theorem, in this light, is
exactly the statement that a symplectomorphism from a neighborhood
of any point to the flat model $(\mathbb R^{2n}, \sum dq^i\wedge
dp_i)$ always exists: "finding canonical coordinates" and "finding a
symplectomorphism to the standard model" are the same task, phrased
in two vocabularies.

\paragraph{Geodesics, connections, and the return of the Christoffel symbols}
The same reinterpretation applies to the free-particle example. Motion
generated by a Riemannian metric --- whether the plain kinetic-energy metric of
a free system, or the Jacobi metric that reformulates a potential as pure
geometry --- is intrinsically \emph{geodesic} motion: $\nabla_{\dot x}\dot x =
0$, where $\nabla$ is the Levi-Civita connection of the metric, an object
defined without reference to any coordinate chart. Written out in coordinates,
this intrinsic equation reads
\begin{equation}
  \ddot x^i + \Gamma^i_{jk}(x)\,\dot x^j \dot x^k \;=\; 0,
  \label{geodesic-coord}
\end{equation}
with exactly the Christoffel symbols that appeared, unbidden, as
"fictitious centrifugal force" in equation~\eqref{polar-free}.
That is not a coincidence or an analogy: a free particle moving in a
straight line, described in polar coordinates, \emph{is} geodesic
motion in the flat metric $ds^2 = dr^2 + r^2 d\phi^2$, and the
Christoffel symbols of that metric in polar coordinates are precisely
the coefficients that turned up as a spurious force term. "Fictitious
force" and "Christoffel symbol" are the same object seen from two
sides. One calls it a force only for as long as one has declined to
think geometrically about it.

\subsection{What Each Form Buys You}

The geometric form earns its keep by telling you, in advance and for
free, which features of a set of equations of motion are physical
and which are artifacts of a chart. It explains why canonical
transformations exist and behave the way they do (Darboux's
theorem). It makes the relationship between symmetry and conservation
essentially automatic: a continuous symmetry is the flow of a vector
field preserving both $\omega$ and $H$, and when that symmetry comes
from the action of a Lie group, the associated conserved quantity is
the group's \emph{momentum map} --- the modern, coordinate-free
statement of Noether's theorem.
It also exposes genuinely global, topological facts that no single
coordinate patch can see: the phase space of a rigid body, for
instance, is the cotangent bundle of the rotation group $SO(3)$,
which cannot be covered by a single nonsingular coordinate chart at
all --- the root cause of gimbal lock, which textbooks often present
as an engineering nuisance but which is, geometrically, an
unavoidable topological obstruction. A coordinate singularity, on
this view, is never a physical singularity: it is only ever the
failure of one particular chart, repaired by changing coordinates or,
better, by working with the coordinate-free structure directly, where
the question of a "singularity" never even arises.

The coordinate form earns its keep by being the only form in which a
problem is actually solved, in closed form or on a computer, and the
only form that connects directly to a measurement, which is always
reported as a number relative to some frame. Darboux's theorem
guarantees good coordinates exist near any point; it says nothing
whatsoever about how to find the ones adapted to a \emph{specific}
Hamiltonian's structure --- action-angle variables for an integrable
system, say, which is frequently the entire technical content of
actually solving a hard problem in mechanics. No amount of abstract
symplectic geometry does that work for you.

Seen from the perspective of the wider project of this book, the two
forms sit on opposite sides of its recurring divide: a coordinate
expression is what gets compared, number for number, against
measurement, and so belongs to the "real world" side of the ledger;
the manifold, the connection, the symplectic form are pieces of
"ideal world" mathematics, whose entire justification is that they
correctly organize and predict the coordinate-level facts, and
nothing more. Neither side is dispensable, and neither is more
fundamental than the other in any sense that would let one be
discarded.

A single example run in both languages makes the trade-off concrete.
A pendulum of length $\ell$ and bob mass $m$, free to swing in any
direction, has configuration space the sphere of radius $\ell$
centered on the pivot.

\emph{Coordinate form.} Using the polar angle $\theta$ from the
downward vertical and the azimuth $\phi$, the Lagrangian is
\begin{equation}
  L \;=\; \tfrac12 m\ell^2\bigl(\dot\theta^2 + \sin^2\theta\,
  \dot\phi^2\bigr) \;+\; mg\ell\cos\theta .
  \label{sph-pendulum-L}
\end{equation}
The coordinate $\phi$ does not appear in $L$ itself, so its conjugate
momentum
\begin{equation}
  p_\phi \;=\; m\ell^2\sin^2\theta\,\dot\phi
  \label{sph-pendulum-pphi}
\end{equation}
is conserved --- it is the angular momentum about the vertical axis
--- and the problem reduces to an effective one-dimensional motion in
$\theta$ alone, with effective potential $V_{\mathrm{eff}}(\theta) =
mg\ell\cos\theta + p_\phi^2/(2m\ell^2\sin^2\theta)$. This is the form
in which the problem is actually solved, e.g.\ by quadrature or
numerically.

\emph{Geometric form.} Configuration space is the sphere $S^2$
itself, with the circle group $SO(2)$ acting on it by rotation about
the vertical axis, an action that preserves the Lagrangian. The
conservation of $p_\phi$ is nothing but the momentum map for this
$SO(2)$ symmetry, an instance of Noether's theorem requiring no
mention of $\theta,\phi$ at all. And the coordinate singularities of
$(\theta,\phi)$ at the poles $\theta=0,\pi$ --- where $\phi$ becomes
undefined and equation~\eqref{sph-pendulum-pphi} degenerates ---
are transparently artifacts of this particular chart: $S^2$ has no
distinguished points and no singularity there whatsoever. A pendulum
hanging straight down, or balanced straight up, is a perfectly smooth
physical configuration; it is only the coordinate system built from
$(\theta,\phi)$ that breaks down at exactly those two points, for the
same reason that all spherical coordinate systems break down at their
poles.

The two languages are not competitors; they specialize in different
questions. The geometric form is the place to ask what is really
happening, structurally, and what is guaranteed to remain true no
matter which coordinates one happens to choose. The coordinate form
is the place --- indeed the only place --- to ask what number a
system actually produces. A physicist who works only in coordinates
is at permanent risk of mistaking an artifact of a chosen chart for a
physical effect, exactly as happened, for a good part of the
twentieth century, with the coordinate singularity of the
Schwarzschild solution at its horizon: a singularity of the
particular $(t,r,\theta,\phi)$ coordinates originally used to write
the metric down, mistaken for decades for a genuine physical
singularity of spacetime itself, and only cleared up by passing to
coordinates (Eddington--Finkelstein, Kruskal) in which the horizon is
manifestly perfectly smooth --- a general-relativistic descendant, on
a much larger stage, of exactly the confusion embodied in the
centrifugal term of equation~\eqref{polar-free}. A physicist who
works only in the geometric language, on the other hand, has an
elegant theory of everything and the means to compute nothing. The
working practice of mechanics has always required both, spoken
fluently, and switched between as the problem at hand demands.
